\chapter{Diagonal Linear Networks and the Lasso Regularization Path}

\section{Introduction}

This chapter formalizes the connections between the training trajectory of diagonal linear networks (DLNs) and the lasso regularization path.
The theoretical background relies on early stopping of gradient flows in neural networks to achieve a tradeoff between sparsity and data fitting.

\section{The \texorpdfstring{$u \circ v$}{u o v} case}

We start by defining the quadratic loss and the related lasso objective. 
Let $M \in \mathbb{R}^{d \times d}$ be a positive semidefinite matrix, and let $r \in \mathbb{R}^d$.

\begin{definition}[Quadratic Loss]\label{def:quadraticLoss}
  \lean{Lasso.quadraticLoss}
  The \emph{quadratic loss} function is defined as
  \[
    \ell(x) = \frac{1}{2} \langle x, M x \rangle - \langle r, x \rangle.
  \]
\end{definition}

\begin{definition}[Lasso Objective]\label{def:lassoObjective}
  \lean{Lasso.lassoObjective}
  For a regularization parameter $\lambda \ge 0$ and an inverse implicit regularization parameter $\mu > 0$, the \emph{lasso objective} is defined as
  \[
    \operatorname{Lasso}(x, \mu) = \ell(x) + \left(\lambda + \frac{1}{\mu}\right) \|x\|_1.
  \]
\end{definition}

\begin{definition}[DLN Objective]\label{def:dlnObjective}
  \lean{Lasso.dlnObjective}
  For weights $u, v \in \mathbb{R}^d$, the \emph{DLN objective} incorporates an explicit weight decay penalty parameterized by $\lambda \ge 0$:
  \[
    L(u, v) = \ell(u \circ v) + \frac{\lambda}{2} \left(\|u\|_2^2 + \|v\|_2^2\right).
  \]
\end{definition}

\begin{definition}[DLN Gradient Flow]\label{def:dlnGradientFlow}
  \lean{Lasso.dlnGradientFlow}
  The gradient flow dynamics of the weights are given by the coupled system of ordinary differential equations
  \[
    \frac{\mathrm{d}u}{\mathrm{d}t} = - \nabla_u L(u, v), \qquad \frac{\mathrm{d}v}{\mathrm{d}t} = - \nabla_v L(u, v),
  \]
  with initializations $u(0) = \sqrt{\varepsilon}\beta$ and $v(0) = \sqrt{\varepsilon}\gamma$.
\end{definition}

\begin{definition}[Average Trajectory]\label{def:averageTrajectory}
  \lean{Lasso.averageTrajectory}
  Let $x^{\varepsilon}(t) = u(t) \circ v(t)$ be the effective linear parameter under the gradient flow. Its time average is
  \[
    \bar{x}^{\varepsilon}(t) = \frac{1}{t} \int_0^t x^{\varepsilon}(t') \, \mathrm{d}t'.
  \]
\end{definition}

\begin{theorem}[Exact Connection to the Lasso]\label{thm:lassoConnectionMonotone}
  \uses{def:lassoObjective, def:dlnGradientFlow, def:averageTrajectory}
  \lean{Lasso.lasso_connection_monotone}
  Assume that the mapping $\mu \mapsto \mu x(\mu)$ is monotonically non-decreasing coordinate-wise. Then for any $s > 0$, as $\varepsilon \to 0$ and tracking the rescaled time $t(s, \varepsilon) = \frac{s}{2} \log(1/\varepsilon)$,
  \[
    \lim_{\varepsilon \to 0} \operatorname{Lasso}(\bar{x}^{\varepsilon}(s), s) = \min_{x \in \mathbb{R}^d} \operatorname{Lasso}(x, s).
  \]
\end{theorem}

\begin{theorem}[Approximate Connection to the Lasso]\label{thm:lassoConnectionApprox}
  \uses{def:lassoObjective, def:dlnGradientFlow, def:averageTrajectory}
  \lean{Lasso.lasso_connection_approx}
  Without the monotonicity assumption, there exists a constant $C > 0$ such that for any $s > 0$,
  \[
    \limsup_{\varepsilon \to 0} \operatorname{Lasso}(\bar{x}^{\varepsilon}(s), s) \le \min_{x \in \mathbb{R}^d} \operatorname{Lasso}(x, s) + C \eta(\lambda, s, z^{\downarrow}(s)),
  \]
  where $z^{\downarrow}(s)$ measures the deviation from monotonicity of $z(\mu) = \mu x(\mu)$.
\end{theorem}

\section{The \texorpdfstring{$u \circ u$}{u o u} case (Positive Lasso)}

We now outline the $u \circ u$ case, which is algebraically simpler and serves as the structural foundation for the $u \circ v$ proofs.

\begin{definition}[Positive DLN Objective]\label{def:posDlnObjective}
  \lean{Lasso.posDlnObjective}
  For a weight vector $u \in \mathbb{R}^d$, the positive DLN objective is:
  \[
    L(u) = \ell(u \circ u) + \lambda \|u\|_2^2.
  \]
\end{definition}

\begin{definition}[Positive DLN Gradient Flow]\label{def:posDlnGradientFlow}
  \lean{Lasso.posDlnGradientFlow}
  The gradient flow dynamics of the weight $u$ is given by
  \[
    \frac{\mathrm{d}u}{\mathrm{d}t} = - \nabla L(u),
  \]
  with initialization $u(0) = \sqrt{\varepsilon}\beta$.
\end{definition}

For this case, the rescaled time is defined as $s = \frac{4}{\log(1/\varepsilon)} t$.

\begin{theorem}[Exact Connection to the Positive Lasso]\label{thm:posLassoConnectionMonotone}
  \uses{def:lassoObjective, def:posDlnGradientFlow}
  \lean{Lasso.pos_lasso_connection_monotone}
  Assume that the mapping $\mu \mapsto \mu x(\mu)$ is monotonically non-decreasing. Then for all $s > 0$,
  \[
    \lim_{\varepsilon \to 0} \operatorname{Lasso}(\bar{x}^{\varepsilon}(s), s) = \min_{x \ge 0} \operatorname{Lasso}(x, s).
  \]
\end{theorem}

\begin{theorem}[Approximate Connection to the Positive Lasso]\label{thm:posLassoConnectionApprox}
  \uses{def:lassoObjective, def:posDlnGradientFlow}
  \lean{Lasso.pos_lasso_connection_approx}
  Without the monotonicity assumption, there exists a constant $C > 0$ such that for any $s > 0$,
  \[
    \limsup_{\varepsilon \to 0} \operatorname{Lasso}(\bar{x}^{\varepsilon}(s), s) \le \min_{x \ge 0} \operatorname{Lasso}(x, s) + C \eta(\lambda, s, z^{\downarrow}(s)).
  \]
\end{theorem}

\subsection{Proof Strategy: Mirror Flow and LCPs}

The proofs of these theorems rely on connecting the DLN dynamics to the positive lasso through two intermediate structures:
\begin{enumerate}
    \item \textbf{Mirror Flow:} The gradient flow can be interpreted as a mirror flow with an entropy potential function, naturally enforcing the positivity constraint.
    \item \textbf{Linear Complementarity Problems (LCPs):} The positive lasso regularization path translates into a parametric LCP, bridging the gap between the primal-dual formulations.
\end{enumerate}
By establishing these connections, Theorem \ref{thm:lassoConnectionMonotone} and \ref{thm:lassoConnectionApprox} (the $u \circ v$ case) follow as corollaries via the structural reduction $p^{\text{pos}} = (u+v)/2$ and $p^{\text{neg}} = (u-v)/2$.
